\section{Πρόβλημα 1}

\subsection{Ερώτηση (i)}

Μας δίνεται το εξής παίγνιο:
\begin{center}
\begin{tabular}{ccccc}
        & $W$ & $X$ & $Y$ & $Z$ \\
    $A$ & 10, 10 & 5, 12 & 0, 15 & 0, 10 \\
    $B$ & 12, 5  & 8, 8  & 4, 10 & 2, 5  \\
    $C$ & 15, 0  & 10, 4 & 6, 6  & 3, 4  \\
    $D$ & 10, 0  & 5, 2  & 4, 3  & 1, 1  \\
\end{tabular}
\end{center}
όπου Θα εκτελέσουμε επαναλαμβανόμενη αφαίρεση ισχυρά κυριαρχούμενων στρατηγικών
και θα δείξουμε ποιες στρατηγικές επιβιώνουν όταν ολοκληρωθεί ο αλγόριθμος.

Σύμφωνα με τον ορισμό, μια στρατηγική $\textbf{p}_i$ για τον παίκτη $i$
κυριαρχεί ισχυρά μια άλλη στρατηγική $\textbf{p}'$ αν
$u_i(\textbf{p}_i, \textbf{p}_{-i}) > u_i(\textbf{p}', \textbf{p}_{-i})$ για
όλες τις διαφορετικές στρατηγικές $\textbf{p}_{-i}$ των άλλων παικτών.

Αρχικά, η στρατηγική $A$ του παίκτη γραμμών κυριαρχείται ισχυρά από την
στρατηγική $C$ ($15 > 10, 10 > 5, 6 > 0, 3 > 0$). Άρα η στρατηγική $A$
αφαιρείται. Παρόμοια, κυριαρχούνται ισχυρά και οι στρατηγικές $B$ και $D$ από την
$C$. Αντίστοιχα, για τον παίκτη στηλών, όλες οι στρατηγικές κυριαρχούνται ισχυρά
από την στρατηγική $X$. Τελικά, οι στρατηγικές που επιβιώνουν είναι οι $C$ και
$X$.

\subsection{Ερώτηση (ii)}

Μας δίνεται το εξής παίγνιο:
\begin{center}
\begin{tabular}{ccccc}
        & $W$  & $X$  & $Y$  \\
    $A$ & 1, 1 & 1, 1 & 0, 0 \\
    $B$ & 1, 1 & 2, 2 & 2, 2 \\
    $C$ & 0, 0 & 2, 2 & 2, 2 \\
\end{tabular}
\end{center}
όπου θα εκτελέσουμε 2 διαφορετικές εκτελέσεις της επαναλαμβανόμενης αφαίρεσης
ασθενώς κυριαρχούμενων στρατηγικών και θα δείξουμε ποιες (διαφορετικές)
στρατηγικές επιβιώνουν όταν ολοκληρωθεί ο αλγόριθμος.

Σύμφωνα με τον ορισμό, μια στρατηγική $\textbf{p}_i$ για τον παίκτη $i$
κυριαρχεί ασθενώς μια άλλη στρατηγική $\textbf{p}'$ αν για τις διαφορετικές
στρατηγικές $\textbf{p}_{-i}$ των άλλων παικτών ισχύει ότι:
$u_i(\textbf{p}_i, \textbf{p}_{-i}) \geq u_i(\textbf{p}', \textbf{p}_{-i})$ και για
τουλαχίστον μία διαφορετική στρατηγική $\textbf{p}_{-i}$ ισχύει ότι:
$u_i(\textbf{p}_i, \textbf{p}_{-i}) > u_i(\textbf{p}', \textbf{p}_{-i})$.

\subsubsection{Εκτέλεση 1}
Η στρατηγική $B$ του παίκτης γραμμών κυριαρχεί ασθενώς την στρατηγική $A$
($1 \geq 1, 2 > 1, 2 > 0$). Παρόμοια κυριαρχείται ασθενώς και η στρατηγική $C$
($1 > 0, 2 \geq 2, 2 \geq 2$). Αυτό μας δίνει το παίγνιο:
\begin{center}
\begin{tabular}{ccccc}
        & $W$  & $X$  & $Y$ \\
    $B$ & 1, 1 & 2, 2 & 2, 2
\end{tabular}
\end{center}

Η στρατηγική $W$ κυριαρχείται ισχυρά και θα μπορούσε να αφαιρεθεί.

\subsubsection{Εκτέλεση 2}
Η στρατηγική $Χ$ του παίκτη στηλών κυριαρχεί ασθενώς την στρατηγική $W$
($1 \geq 1, 2 > 1, 2 > 0$). Παρόμοια κυριαρχείται ασθενώς και η στρατηγική $Y$
($1 > 0, 2 \geq 2, 2 \geq 2$). Αυτό μας δίνει το παίγνιο:
\begin{center}
\begin{tabular}{ccccc}
        & $X$  \\
    $A$ & 1, 1 \\
    $B$ & 2, 2 \\
    $C$ & 2, 2 \\
\end{tabular}
\end{center}

Η στρατηγική $A$ κυριαρχείται ισχυρά και θα μπορούσε να αφαιρεθεί.
